<<<P001395> Тип поощрения
<<P001385>

<<<P001294>%
<<<P001396>{ifElement} взаимодействует с продвижением типа
Таким же образом, как <<<P001397> Заявления делают.
<<<P000006> быть <<<P001398>{ifElement} формы
<<P001043> или
<<P001044>.
Если <<P000012> показывает локальную переменную <<<P000013> имеет тип <<<P000014>
Тип <<<P000015> Известен как <<<P000016> в <<<P000017>,
Если ни одно из следующих утверждений не верно:

<<P000899>
<<P001399> <<P000018> потенциально мутирует в <<<P000019>,
<<P001400> <<<P000020> Потенциально мутирует в функции
кроме той, где <<<P000021> объявляется, или
<<<P001401> <<P000022> Доступ осуществляется с помощью функции, определенной в <<<P000023> и
<<P000024> Потенциально мутирует в любом месте в пределах <<<P000025>.
<<P000929>

% % TODO(Eernst): Приходите, обновите это.
<<<P001402>{%>
Продвижение типа, вероятно, станет более сложным в будущей версии Dart.
Когда это происходит,
<<<<P001403>> Элементы будут продолжать соответствовать <<<P001404> заявления
(<<P000959>)%
?


<<<P001405> Литературная оценка элементов
<<P001386>

<<<P001295>%
Оценка последовательности сбора буквальных элементов
(<<P000960>)
доходность a
\IndexCustom{сбор буквальной последовательности объектов}{%
Коллекция буквально! объектная последовательность.
Также называется <<<P001407>{объектная последовательность}, когда не может возникнуть двусмысленности.

<<<P001296>%
Мы используем нотацию
% <<<P001408> Не может быть использован после '@', поэтому мы используем приближение.
\IndexCustom{\LiteralSequence{\ldots}}{[[...]]@%
<<<<P001412>< \hspace{-0.6mm}[<<P001414>>><<P001415>>>{-0.6mm}]}}
обозначать объектную последовательность с явно перечисленными элементами,
и мы используем '<<<P000026>' для вычисления сцепления последовательностей объектов
\commentary{(как в $s_1 + s_2$)},
Это операция, которая будет успешной и не будет иметь побочных эффектов.
Каждый элемент в последовательности является объектом. <<P000028> или парой <<<P000029>.
Не существует понятия типа элемента для последовательности объекта,
Не существует динамических ошибок, возникающих из
Тип несоответствия во время конкатенации.

<<<P001417>{%>
Последовательности объектов можно безопасно рассматривать как механизм низкого уровня.
которые могут пропустить другие необходимые действия, такие как динамические проверки типа
Каждый доступ к объектной последовательности происходит
в коде, созданном языком, определяемым desugaring
на статически проверенных конструкциях.
Остается неустановленным, как реализуется объектная последовательность,
Требуется только, чтобы он содержал указанные объекты или пары.
в данном порядке.
Для каждой коллекции,
последовательность используется в данном порядке для заполнения коллекции,
в том виде, который специфичен для данного вида,
который указывается отдельно
(<<P000961>, <<P000962>, <<<P000963>>%

Может существовать реальная структура данных.
представление последовательности объекта во время выполнения,
но последовательность объектов также может быть устранена, например,
потому что каждый элемент вставляется непосредственно в целевой набор
Как только она будет вычислена.
Обратите внимание, что каждая последовательность объектов будет содержать исключительно объекты.
Или же он будет содержать только пары,
Любая попытка создать смешанную последовательность вызовет
Ошибка во время компиляции или во время выполнения
(последний может иметь место для элемента с разбросом со статическим типом \DYNAMIC{}.%)
?

<<<P001297>%
Допустим, что дается дословный сборник \metavar{target},
и будет использоваться последовательность объектов, полученная, как описано ниже
населять <<<P001420>{цель}.
<<<P000030> обозначает динамический тип <<<P001421>{цель}.

<<<P001422>{%>
Доступ к типу <<<P001423>{цель} необходим ниже
Для увеличения динамических ошибок в определенных точках во время
Оценка объектной последовательности.
Обратите внимание, что динамический тип <<<P001424>{цель} статически известен.
за исключением связывания переменных любого типа в его <<<P001425>{тип Аргументы}.Это означает, что на некоторые вопросы можно ответить во время компиляции, например:
Или нет (P001045). происходит как суперинтерфейс
<<P000031>.
В любом случае, <<<P000032> гарантируется осуществление
<<P001046> (когда <<<P001426>{цель} - список или набор)
или <<<P001047> (когда <<<P001427>{цель} является картой),
но никогда оба.%
?

<<<P001298>%
Предположим, что дается место в коде и динамический контекст,
Таким образом, возможна обычная оценка выражения.
<<<P001428> Оценка последовательности буквальных элементов коллекции {%}
Буквальный элемент коллекции! оценка последовательности
В этой связи следует отметить, что в этой связи следует принять во внимание, что в данном случае

<<<P001299>%
Пусть $s_{\metavar{syntax}}$ формы \List{\ell}{1}{k} быть
Последовательность сбора буквальных элементов.
Последовательность объектов <<<P000034>
полученных путем оценки <<P000035>
является сцепление последовательностей объектов
Полученные путем оценки каждого элемента <<P000036>, <<P000037>:
<<P000038>,
где \EvaluateElement{<<P001432>>>_j} обозначает объектную последовательность, полученную
Оценка одного литературного элемента коллекции <<P000039>.

<<<P001300>%
При псевдозаявлении формы
<<<P001048>
используется в нормативном кодексе ниже,
обозначает расширение <<<P000041> с последовательностью объекта
по результатам оценки <<<P000042>,
но также обозначает спецификацию действий, предпринятых для
Последовательность объектов,
и для получения побочных эффектов, связанных с этим вычислением,
как подразумевается при оценке выражений и исполнении заявлений
Как указано ниже для оценки \EvaluateElement{\ell}.

<<<P001301>%
При псевдозаявлении формы
<<P001049>
происходит в нормативном кодексе ниже,
Это происходит в точке, где вычисление завершено.
и указывает, что значение \EvaluateElement{\ell} составляет <<<P000044>.
<<<P001437> Оценка литературного элемента коллекции <<<P000045>{%}
Буквальный элемент коллекции! оценка
в данном контексте
к объектной последовательности
\IndexCustom{\EvaluateElement{\ell>>>}}
Оценка Element(l)@\emph Элемент\code{(<<<P000046>>>)}
Затем указывается следующее:

<<<P001302>%
<<<P001442> Элемент выражения
В данном случае <<<P000047> выражение <<<P000048>;
<<P000049> оценивается на объект <<<P000050>
<<P000051>.
<<P001443>

<<<P001303>%
<<<P001444> Элемент карты
В данном случае <<<P000052> Это пара выражений
<<<P001051>;
<<<P000055> оценивается на объект <<<P000056>,
<<<P000057> оценивается на объект <<<P000058>,
<<P000059>.
<<P001445>

<<<P001304>%
<<<P001446> Элемент распространения.
Элемент <<<P000060> Он имеет форму <<<P001052> или <<<P001053>.
Оценить <<<P000063> на объект <<<P000064>.

<<P000900>
<<P001447>
Когда <<<P000065> означает <<<P001054>:
% % TODO(Eernst): Приходите в NNBD, эта ошибка больше не может произойти: удалите.
Если <<<P000067> Это нулевой объект, после чего происходит динамическая ошибка.
В противном случае оценка осуществляется на этапе 2.

Когда <<<P000068> означает <<<P001055>:
Если <<<P000070> Тогда нулевой объект
<<P000071>.
В противном случае оценка осуществляется на этапе 2.
<<P001448>
<<<P000072> Динамический тип <<<P000073>.
<<<P000074> Статический тип <<<P000075>.
Когда <<P000076> Не является высшим типом
(<<P000964>),
Пусть <<<P000077> будет <<<P000078>.
Когда <<<p000079> Это топовый тип:
Если <<<P001449>{цель} - это список или набор
Пусть <<<P000080> будет <<<P001056>;
иначе
\commentary{(где \metavar{цель} - карта)},
Пусть <<<P000081> будет <<<P001057>.

<<P000901>
<<P001452>
Когда <<<P001453>{цель} - список или набор
<<P000082> реализация
(<<P000965>)
<<<P001058>,
следующий код выполняется в контексте, где <<<P000083> происходит,
где <<<P001059>, <<P000084>, <<<P001060> и <<<P001061> свежие переменные,<<<P001062> является переменной нового типа, связанной с
Аргумент фактического типа <<<P000085> в <<<P001063>
(<<P000966>):

\vspace{-2ex}<<P000902>>>[t]{<<P001455>>>]
<<P000003>
<<P000930>

<<<P001456>{%>
Код использует псевдовариабельный <<<P000086> обозначает последовательность объектов.
Мы не указываем тип <<<P000087>,
Эта переменная используется только для указания требуемого значения.
Семантические действия, предпринятые для сбора полученной объектной последовательности.
В случае, если реализация не
Изображение <<<P000088> Во всяком случае,
Действие может заключаться в немедленном продлении <<<P001457>{target}.
Аналогичный подход используется в последующих случаях.%
?
<<<P001458>
Когда <<<P001459>{цель} - карта и <<P000089> реализация
<<<P001064>,
следующий код выполняется в контексте, где <<<P000090> происходит,
где <<<P001065>, <<P000091>, <<<P001066> <<<P001067> и <<<P001068>
свежие переменные,
и <<<P001069> и <<<P001070> свежие переменные типа, связанные с
первый, соответственно второй фактический тип аргумента
из <<<P000092> по адресу <<<P001071>:

\vspace{-2ex}<<P000903>>>[t]{<<P001461>>>]
<<P000004>
<<P000931>

% не изменится при nnbd: аргументы типа «распространение» могут быть <<<P001462>.
Реализация может задерживать динамические ошибки
Это происходит, если данный <<P001072> не имеет типа <<<P001073>,
или отданное <<P001074> не имеет своего типа <<<P001075>,
Но это не может произойти после того, как пара была добавлена к <<P000093>.
<<P001463>
В противном случае возникает динамическая ошибка.

<<<P001464>{%>
Это происходит, когда цель является повторяемой соответственно картой.
Распространения нет, что возможно для
спред, статический тип которого <<<P001465>.%
?
<<P000932>
<<P000933>

<<<P001466>{%>
Это может быть не самый эффективный способ перемещения предметов в коллекции.
Реализация может, конечно, использовать любой другой подход.
с таким же наблюдаемым поведением.
Тем не менее, чтобы дать реализациям больше возможностей для оптимизации.
Допустим также следующее.%
?

<<<P001305>%
Если <<P000094> Объект, динамический тип которого реализует
(<<P000967>)
<<<P001076>, <<<P001077> или <<<P001078>,
Реализация может выбрать вызов <<<P001079> на объекте.
Если <<P000095> является объектом
чей динамический тип <<<P001080>,
реализация может выбрать вызов оператора <<<P001467>
Для доступа к элементам из списка.
Если он это сделает, он будет только передавать индексы.
которые являются неотрицательными и меньше значения, возвращенного <<<P001081>.

<<<P001468>{%>
Это позволяет использовать более эффективный код для
Распределение коллекции и доступ к ее частям.
Эти классы должны иметь
Эффективная и бесперебойная реализация
<<<P001082>> и оператора <<<P001469>{[]}.
Реализация Dart может определить, применяются ли эти варианты.
в момент компиляции на основе статического типа <<<P000096>,
или во время выполнения на основе фактического значения.%
?
<<P001470>

<<<P001306>%
<<<P001471> Если элемент
Когда <<<P000097> <<<P001472>{ifElement} формы
<<P001083> или
<<<P001084>,
Состояние <<<P000103> оценивается по значению <<<P000104>.
Если <<<P000105> <<<P001473> затем
<<P000106>.
Если <<<P000107> <<<P001474>{} и <<<P000108> Присутствует тогда
<<P000109>,
Если <<<P000110> Не присутствует тогда
<<P000111>.
% <<<P000112> может иметь тип <<<P001475>.
Если <<P000113> не является ни \TRUE{}, ни \FALSE{} Затем происходит динамическая ошибка.
<<P001478>

<<<P001307>%
<<<P001479> Для элемента
<<<P000114> быть получено из <<<P001480>{<forLoopParts>} и
Пусть <<<P000115> будет <<<P001481>{для элемента} формы
<<<P001085>,Где «P001482»? Это означает, что <<<P001483> может присутствовать или отсутствовать.
Для оценки <<P000118>,
следующий код выполняется в контексте, где <<<P000119> происходит,
где <<<<P001484> Присутствует тогда и только тогда, когда присутствует в <<<P000120>:

\vspace{-2ex}<<P000904>>>[t]{<<P001486>>>]
<<P000005>
<<P000934>
<<P001487>


<<<<P001488>{Список литературы}
<<P001387>

<<<P001308>%
В этом разделе указывается, как буквальный список \metavar{list} пройден и
<<<P001490>{предполагаемый тип элемента}{список буквальный! Тип элемента
Для <<<P001491>{список} определяется.
Сначала мы указываем, как определить тип элемента одного элемента.
Как использовать этот результат для вывода
Тип элемента <<<P001492>{список} в целом.

<<<P001309>%
Тип контекста <<<P000121>
(<<P000968>)
Для каждого элемента <<<P001493>
Получено из контекстного типа <<<P001494>{список}.
Если вывод вниз ограничивает тип <<<P001495>
<<<P001086> или <<P001087> Для некоторых <<P000124>
<<<P000125> <<<P000126>.
В противном случае <<<P000127> \FreeContext{}
(P000969).

<<<P001310>%
<<<P000128> Термин, полученный из <<<P001497>{элемент}.
Вывод типа элемента <<<P000129> Тип контекста <<<P000130>
Происходит следующим образом,
где тип контекста для вывода типа элемента всегда <<<P000131>,
Если ничего не сказано об обратном:

<<<P001311>%
<<<P001498> Элемент выражения
В данном случае <<<P000132> Выражение <<<P000133>.
Тип элемента <<<P000134> это
% Тип <<<P000135> Не является типом элемента, поэтому мы упоминаем <<P000136>.
Предполагаемый тип <<<P000137> в контексте <<P000138>.
<<P001499>

<<<P001312>%
<<<P001500> Элемент карты
% Мы не можем точно сказать, является ли это ошибкой в <<P000970> возникает первым
% (скажем, когда буквальный список имеет явный тип аргумента, он не является
% гарантирует, что реализация выполняет вывод типа на ней вообще.
% Таким образом, мы _must_ укажем ошибку в <<<P000971>, и это местоположение должно
% — это комментарий, потому что у нас уже есть ошибка в другом месте.
<<<P001501>{%>
Этого не может произойти:
Это ошибка времени компиляции
Когда листовой элемент списка буквально является элементом карты
(<<P000972>)%
?
<<<P001502>

<<<P001313>%
<<<P001503> Элемент распространения.
<<<P000139> Выражение <<P000140>.
Если <<<P000141> означает <<<P001088>,
<<<P000143> Выберите тип <<<P000144> в контексте <<<P001089>.
Иначе
<<<P001504>{(когда <<<P000146> является <<<P001090>>)},
%% TODO (eernst): Приходите в NNBD, добавьте <<<P000973> в спецификацию
% — «неотменяемый тип».
<<<P000148> быть ненулевым типом
% % TODO(Eernst): Уточнить, будет ли вывод действительно иметь такой контекст;
% ясно, что мы должны устранить «Null» из <<<P000149>, а также что $e$
% допускается иметь потенциально нулевой тип, и это кажется неудобным.
% %, если мы используем 'Iterable<<<<P000151>> в качестве контекстного типа и отказа в случае, когда <<<P000152>,
% % говорят: «Список <<<<P000153>>?» Для некоторых <<P000154>.
Предполагаемый тип <<<P000155> в контексте <<<P001091>.

<<P000905>
<<<P001505>
Если <<P000157> реализация <<<P001092>,
Тип элемента <<<P000158> это
Аргумент типа <<<P000159> при <<<P001093>.
<<<P001506>
Если <<<P000160> <<<P001507>,
Предполагаемый тип элемента <<<P000161> является <<<P001508>.
<<P001509>
Если <<<P000162> <<<P001094> и оператор спреда <<<P001510>{...?},
Тип элемента <<<P000163> является <<<P001095>.
<<P001511>
В противном случае возникает ошибка времени компиляции.
<<P000935>
\vspace{-5 мм)
<<<P001513>

<<<P001314>%
<<<P001514> Если элемент
В данном случае <<<P000164> имеет форму
<<P001096> или
<<P001097>.
Состояние <<<P000170> Всегда выводится контекстный тип <<<P001098>.
Допустим, что "<<<P001099>" отсутствует.
Затем, если указанный тип элемента <<<P000172> является <<<P000173>,
Тип элемента <<<P000174> является <<<P000175>.

В противном случае присутствует «\code{\ELSE\,\,<<<P000176>>>}».
Если выбранный тип элемента <<<P000177> является <<<P000178> и
Тип элемента <<<P000179> является <<<P000180>,
Тип элемента <<<P000181> это
наименьшая верхняя граница <<<P000182> и <<<P000183>.
<<P001515>

<<<P001315>%
<<<P001516> Для элемента
В данном случае <<<P000184> имеет форму
<<P001101>
где <<<P000187> является производным от <<<P001517>{forLoopParts} и
<<<P001518>> Это означает, что <<<P001519> может присутствовать или отсутствовать.

Ошибки времени компиляции для <<<P000188> как
Ошибки, которые могут возникнуть с соответствующим <<<P001520> заявление
<<<P001521><<<P001522> \,\,\FOR\,\,$P$<<P002029>>>\,\{\}>,
находится в том же объеме, что и <<<P000190>.
Кроме того, ошибки и анализ типов <<<P000191> выполняется
как если бы это произошло в пределах объема тела указанного <<<P001524>{} заявления.
<<<P001525>{%>
Например, если <<<P000192> имеет форму
<<<p001102>
переменная <<<P001103> Находится в пределах <<<P000194>.%
?

Вывод для частей
<<<P001526>{%>
(таких, как выразительное выражение фор-ин,
или \synt{forInitializerStatement} для петли)%
?
выполнено как для соответствующего \FOR{} заявление,
включая <<<P001529>< если и только если элемент содержит <<<P001530>.
Затем, если указанный тип элемента <<<P000195> является <<<P000196>,
Предполагаемый тип элемента <<<P000197> является <<<P000198>.

<<<P001531>{%>
Другими словами, вывод течет вверх от элемента тела.
?
\vspace{3 мм)
<<P001533>

<<<P001316>%
Наконец, мы определяем
\IndexCustom{тип вывода в списке буквальный}{%
Типовой вывод! список буквальный
в целом.
Предположим, что \metavar{list} является производным от \synt{listLiteral}
и содержит элементы \List{\ell}{1}{n},
и тип контекста для \metavar{list} является <<<P000199>.

<<P000906>
<<P001540>
Если <<<P000200> <<<P001541> затем
Предполагаемый тип элемента для \metavar{list} является <<<P000201>,
где <<P000202> является наименьшей верхней границей
Выводимые типы элементов \List{\ell}{1}{n}.

<<P001545>
% % TODO(Eernst): Спецификация функции говорит <<<P000203>, но как мы знаем <<<P000204> Это тип?
В противном случае
Предполагаемый тип элемента для <<<P001546>{list} является <<<P000205>,
где <<<P000206> определяется путем нисходящего вывода.
<<P000936>

<<<P001317>%
В обоих случаях статический тип <<<P001547>{список} является <<<P001104>.


<<<<P001548>{Список}
<<P001388>

<<<P001318>%
A <<<P001549>{список буквальный}{буквальный}! список
обозначает объект списка, который является целым индексируемым набором объектов.
Правило грамматики для <<<P001550>{listLiteral} указано в другом месте.
(P000974).

<<<P001319>%
При заданном списке буквальный <<<P000208> Нет никаких аргументов,
Тип аргумента <<P000209> Выбирается так, как указано в другом месте
(<<P000975>),
и <<<P000210> отныне рассматривается как
(<<P000976>)
<<P001105>.

<<<P001551>{%>
Статический тип списка буквальный формы <<<P001106>
<<<P001107>
(<<P000977>)%
?

<<<P001320>%
<<<P000216> быть буквальным списком формы
<<<P001552><<<<P000217>> \List{<<P001554>>>}{1}{m}]}.
Это ошибка времени компиляции, если элемент листа <<<P000218> является
<<<P001555>{mapElement}.
Это ошибка времени компиляции, если для некоторых <<<P000219>
<<P000220> не имеет элемента типа,
Тип элемента <<<P000221> Не может быть назначено <<<P000222>.

<<<P001321>%
Список может содержать ноль или более объектов.
Количество объектов в списке – это их размер.
Список имеет связанный набор индексов.Пустой список имеет пустой набор индексов.
Непустый список имеет индекс <<<P000223>
где <<P000224> Это размер списка.
Динамическая ошибка при попытке получить доступ к списку
использовать индекс, который не входит в его набор индексов.

<<<P001556>{%>
Системные библиотеки определяют множество членов для типа <<<P001108>,
Но мы указываем только минимальный набор требований.
Используется самим языком.%
?

<<<P001322>%
Если список буквальный <<<P000225> начинается с зарезервированного слова <<<P001557>
<<<P000226> происходит в постоянном контексте
(<<P000978>),
Это а
<<<P001558>{постоянный список буквально}{буквальный!список!постоянный},
Что является постоянным выражением
(<<P000979>)
и поэтому оценивается во время компиляции.
В противном случае это a
<<<P001559>{список времени выполнения буквальный}{буквальный!список!срок выполнения}
и оценивается во время выполнения.
Только список времени выполнения буквы могут быть мутированы
После того, как они были созданы.
% Эта ошибка может произойти потому, что постоянное является динамическим свойством.
Попытка мутировать постоянный список буквально приведет к динамической ошибке.

<<<P001560>{%>
% Прямо или косвенно верно следующее: Существует <<<P001561>
% модификатор в списке буквальный, или список буквальный в целом
% в постоянном контексте.
Обратите внимание, что в сборнике буквальные элементы постоянного списка буквальные
происходит в постоянном контексте
(<<P000980>),
что означает <<<P001562> Модификаторы не должны быть указаны явно.%
?

<<<P001323>%
Это ошибка времени компиляции
Если элемент постоянного списка буквален, то он не постоянен.
Это ошибка времени компиляции, если аргумент типа постоянного списка буквален.
<<<P001563>{(независимо от того, является ли это явным или предполагаемым)}
Не является постоянным выражением типа
(P000981).

<<<P001564>{%>
Связывание формального параметра типа замкнутого класса или функции
не известны во время составления,
Поэтому мы не можем использовать такие параметры типа в постоянных выражениях.
?

<<<P001324>%
Значение постоянного списка буквально
<<<P001565><<<P001566> <<<P002033><<<P002034><<<<P000227>> \List{<<P001568>>>}{1}{m}}
является объектом <<<P000228> чей класс реализует встроенный класс
<<P001109>
где <<P000230> - фактическое значение <<<P000231>
(<<P000982>),
и содержимое которого представляет собой объектную последовательность <<<P001569>{o}{1}{n}
Оценка \List{<<P001571>>>}{1}{m}
(P000983).
<<<P000232> <<<P000233> (по индексу <<P000234>) В этом случае <<P000235>.

<<<P001325>%
<<<P001110>
<<<P001111>
Будьте двумя постоянными буквами.
Пусть <<<P000240> с содержимым <<<P000241> Реальный тип аргумента <<P000242>
соответственно
<<P000243> Содержание <<<P000244> Реальный тип аргумента <<P000245>
быть результатом их оценки.
Далее <<<P001112> Показатель <<<P001572>{} если {f}
<<<P001113> и <<<P001114>
<<<P001573>{} для всех <<<P000252>.

<<<P001574>{%>
Другими словами, постоянные буквальные списки канонизированы.
Нет необходимости рассматривать канонизацию.
для других экземпляров типа <<<P001115>,
Поскольку такие случаи не могут быть
результат оценки постоянного выражения.%
?

<<<P001326>%
Список времени выполнения буквальный
<<<P001575><<<<P000253>> \List{<<P001577>>>}{1}{m}}
Оценка проводится следующим образом:

<<P000907>
<<P001578>
Оцениваются элементы \List{\ell}{1}{m}
(<<P000984>),
в объектную последовательность \LiteralSequence{\List{o}{1}{n}}.
<<P001583>
Свежий пример (<<P000985>) <<<P000254>, размер <<<P000255>,
чей класс реализует встроенный класс <<<P001116>
выделяется,
где <<P000257> - фактическое значение <<P000258>
(<<P000986>).
<<P001584>
Оператор \lit{[]=} вызывается на $o$ с
Первый аргумент <<P000260> Второй аргумент
<<P000261>.
<<P001586>
Результатом оценки является <<<P000262>.
<<P000937>

<<<P001327>%Объекты, созданные буквами списка, не переопределяются
<<<P001587>{==} оператор, унаследованный от P001117 класс.

<<<P001588>{%>
Обратите внимание, что в этом документе не указывается приказ
в котором установлены элементы.
Это позволяет выполнять параллельные задания в списке.
Если этого желает реализация.
Порядок может соблюдаться только следующим образом (и на него нельзя полагаться):
Если элемент <<<P000263> не является подтипом элемента типа списка,
Ошибка динамического типа возникает, когда <<P000264> присваивается <<<P000265>.%
?


<<<<P001589>{Сет и карта Буквальная неопределенность}
<<P001389>

<<<P001328>%
Некоторые термины, такие как <<<P001590><<<P002035><<<P002036>> и <<<P001591><<<P002037><<<P002038> \id\,\}> являются двусмысленными:
Они могут быть либо буквальными, либо буквальными.
Эта двусмысленность устраняется в два этапа.
На первом этапе используется только тип синтаксиса и контекста.
% % TODO (Eernst): Включите эту ссылку, когда будет определен «тип контекста».
% % (<<P000987>)
и описано в этом разделе.
Второй этап использует типы выражений и описывается следующим образом.
(P000988).

<<<P001329>%
Пусть <<<P000266> будет <<<P001593>{setOrMapLiteral}
с помощью элементов «<<<<<<<<<<<<<<<<<<<<<<<<<<<<<<<<<<<<<<<<<<<<<<<<<<>>>>>>>>>><<<<<<<>>>>>>>>>>>>>>>>>>>>>>>>>>>>>>>>>>>>>>>>>>>>>>>>>>>>>>>>>>>>>>>>>>>>>>>>>> и типами «Публикация»
Если <<<P000269> <<<P001594> Тогда пусть <<<P000270> быть неопределенным.
% % TODO(Eernst): Определить «наибольшее закрытие» типа контекста
% при определении «контекстного типа».
В противном случае пусть <<<P000271> быть наибольшей закрытостью <<<P001595> С)
(P000989).

% % TODO(Eernst): Исключить слова "тип контекста", "наибольшее закрытие".
<<<P001596>{%>
В будущей версии этого документа будут указаны типы контекста.
Основная интуиция заключается в том, что
<<P001039>
Это тип, объявленный для
Получающая организация, такая как
формальный параметр <<P000272> или объявленной переменной <<<P000273>.
Этот тип будет контекстным для
фактический аргумент, переданный <<<P000274>,
соответственно начальное выражение для <<<P000275>.
В некоторых ситуациях контекст не имеет ограничений.
Например, когда переменная объявлена с помощью <<<P001597>{} Вместо аннотации типа.
Это приводит к возникновению
<<<P001598>{неограниченный тип контекста}{контекстный тип!неограниченный},
\IndexCustom{<<P001600>>><<P001601>>>}{[@<<P001602>>>],
Это может также происходить в сложном термине, например, <<<P001118>.
% % TODO(Eernst): Уточните, почему мы не используем i2b
% - введение понятия наибольшего (и наименьшего) закрытия.
Наибольшее закрытие контекстного типа <<<P000276> это
Наименее распространенный супертип всех типов
Получается путем замены \FreeContext{} по типу.%
?

<<<P001330>%
Шаг неоднозначности этого раздела является
первая применимая запись в следующем списке:

<<P000908>
<<<P001604> Когда <<P000277> Типовые аргументы <<<P001605>> T}{1}{k} <<P000278>:
Если <<<P000279>, то <<<P000280> Буквальный набор со статическим типом
<<P001119>.
Если <<<P000282>, то <<<P000283> Буквальная карта со статическим типом
<<P001120>.
В противном случае возникает ошибка времени компиляции.
<<<P001606>
Когда <<P000286> реализация
(<<P000990>)
<<P001121> но не <<<P001122>,
<<P000287> Это набор буквальный.
Когда <<P000288> Реализация <<<P001123> но не <<<P001124>,
<<<P000289> - буквальная карта.
<<<P001607>
Когда <<P000290>
\commentary{(то есть <<P000291>) имеет элементы листьев)}:
Если <<P000292> Содержит <<<P001609>{mapElement}
а также <<<P001610>{выражениеЭлемента},
Происходит ошибка времени компиляции.
В противном случае, если <<<P000293> содержит <<<P001611>{выражениеЭлемента},
<<<P000294> - это буквальный набор.
В противном случае <<P000295> содержит \synt{mapElement} и $e$ Это буквальная карта.
<<<P001613>
Когда <<P000297> имеет форму \code{\{\}} и $S$ неопределенным,
<<<P000299> - буквальная карта.
<<<P001615>{%>Нет более глубокой причины для такого выбора.
Но факт, что \code{<<P002043>>>\}} Это карта по умолчанию
Он был полезен, когда были введены буквы,
Потому что это было бы переломным изменением, чтобы сделать его установленным.
?
<<<P001617>
В противном случае <<<P000300> остается неоднозначным.
<<<P001618>{%>
В данном случае <<<P000301> не является пустым, но содержит только спреды
0 или более раз в <<<P001619> Элементы или <<<P001620>{для элементов}s.
Затем во время вывода будет происходить раздвоение
(<<P000991>)%
?
<<P000938>

<<<P001621>{%>
Когда этот шаг не определяет статический тип,
будет определяться по типу вывода
(<<P000992>)%
?

<<<P001331>%
Если этот процесс успешно дезавуирует буквальный
Мы говорим, что <<P000302> это
<<<P001622>{неоднозначно множество}{set!unambiguously}
или
<<<P001623>{неоднозначно карта}{map!unambiguously},
по мере необходимости.


<\subsubsection{Сет и карта буквального вывода}
<<P001390>

<<<P001332>%
В этом разделе указывается, как <<<P001625> Литературный <<P000303> пройденный
и связанных с ним
<<<P001626>{предполагаемый тип элемента}{набор или карта буквально! Тип элемента
и/или связанный
<<<P001627>{пара типа ключа и значения}{%
Настройка или карта буквально! пара типа ключа и значения
определяется.

<<<P001628>{%>
Если <<P000304> имеет тип элемента, то он может быть набором,
Если у него есть пара типа ключа и значения, то это может быть карта.
%
Однако, если буквальный <<P000305> содержит разветвленный элемент типа <<<P001629>,
Этот элемент не может быть использован для определения <<<P000306> Это набор или карта.
Неоднозначность представлена как имеющая \emph{оба}
Тип элемента и пара типа ключа и значения.

Ошибка, если неясность не определена другими элементами,
Но если это решено, то требуется динамический элемент распространения.
Чтобы оценить подходящий пример
(реализация <<<P001125>, когда <<<P000307>) Это набор,
и реализации <<<P001126>, когда <<<P000308> Это карта,
Это означает, что это динамическая ошибка, если есть несоответствие.
В других ситуациях ошибка компиляции состоит в том, что оба
тип элемента и пара типа ключа и значения,
<<<P000309>> Это должен быть и набор, и карта.
Вот пример: %
?

<<P000000>

<<<P001333>%
Пусть <<<<P001631>{коллекция} будет сборником буквально
полученный из <<<P001632>{setOrMapLiteral}.
Предполагаемый тип <<<P001633>{элемент} является типом элемента <<<P000310>,
пара типа ключа и значения <<<P000311> или оба.
Вычисляется относительно типа контекста <<<P000312>
(<<P000993>),
который определяется следующим образом:

<<P000909>
<<<P001634>
Если <<<P001635>{коллекция} является однозначно установленным
(<<P000994>)
<<<P000313> <<<P001127>,
% % TODO(Eernst): Добавьте ссылку, когда мы укажем вывод.
где <<P000315> определяется нисходящим выводом,
и может быть <<<P001636>
% % TODO(Eernst): Исправьте эту ссылку, когда мы указываем типы контекста.
(<<P000995>)
Если нисходящий вывод не ограничивает его.

% % TODO(Eernst): Удалите это, когда мы укажем вывод.
<<<P001637>{%>
В будущей версии этого документа будет указан вывод,
Понятие нисходящего вывода,
и ограничения.
Краткая интуиция заключается в том, что вывод выбирает значения для
Типовые параметры в генерических конструкциях
где не приводятся аргументы,
нацеливание на тип, который соответствует данному типу контекста;
Нисходящий вывод делает это путем передачи информации.
из данного выражения в его подвыражения,
Вывод вверх распространяет информацию в противоположном направлении.
Ограничения выражаются в терминах контекстных типов;
Быть неограниченным означает иметь <<<P001638>{} в качестве контекстного типа.
Имеет тип контекста, который <<<P001639>
один или несколько случаев <<<P001640>
обеспечивает частичное ограничение на предполагаемый тип.%?
<<<P001641>
Если <<<P001642>{коллекция} является однозначной картой
<<<P000316> <<<P001128>
где <<<P000319> и <<<P000320> определяется нисходящим выводом,
и может быть <<<P001643>
Если нисходящий контекст не ограничивает одно или оба.
<<<P001644>
В противном случае <<<P001645>{коллекция} неоднозначна,
и нисходящий контекст для элементов <<<P001646>
<<<P001647>.
<<P000939>

<<<P001334>%
Мы говорим, что коллекция буквальный элемент
<<<P001648>{может быть множество}{сборник буквальный элемент! может быть набором.
если он имеет тип элемента;
это
<<<P001649>{может быть карта}{сборник буквальный элемент! Это может быть карта
если он имеет пару типа ключа и значения;
это
<<<P001650>{должен быть набор}{сборник буквальный элемент! Должен быть набор.
если она может быть установленной и не иметь пары типа ключа и значения;
и это
<<<P001651>{должна быть карта}{сборник буквальный элемент! Должна быть карта
может быть картой и не имеет типа элемента.

<<<P001335>%
<<<P000321> Термин, полученный из <<<P001652>{элемент}.
<<<P001653> Вывод типа {%}
Типовой вывод! Буквальный элемент коллекции
<<P000322> Тип контекста <<<P000323> После этого доходы
следующим образом:

<<<P001336>%
<<<P001654> Элемент выражения
В данном случае <<P000324> Выражение <<P000325>.
%
Если <<<P000326> <<<P001655>,
Тип элемента <<<P000327> это
Предполагаемый тип <<<P000328> в контексте <<<P001656>.
%
Если <<<P000329> <<<P001129>,
Тип элемента <<<P000331> это
Предполагаемый тип <<<P000332> в контексте <<P000333>.
<<<P001657>

<<<P001337>%
<<<P001658> Элемент карты
В данном случае <<P000334> Это пара выражений <<<P001130>.
%
Если <<<P000337> <<<P001659>,
Выводимая пара типа ключа и значения <<P000338> является <<<P000339>,
где <<<P000340> и <<<P000341> это
Предполагаемый тип <<<P000342> соответственно <<P000343>,
в контексте <<<P001660>.
%
Если <<<P000344> <<<P001131>,
Выводимая пара типа ключа и значения <<P000347> является <<<P000348>,
где <<P000349> это
Предполагаемый тип <<<P000350> в контексте <<P000351> и
<<P000352> это
Вычисленный тип <<<P000353> в контексте <<P000354>.
<<<P001661>

<<<P001338>%
<<<P001662> Элемент распространения.
В данном случае <<P000355> имеет форму
<<<P001132> или <<P001133>.
Если <<<P000358> <<<P001663> <<<P000359>> быть
Предполагаемый тип <<<P000360> в контексте <<<P001664>.
Тогда:

<<P000910>
<<<P001665>
Если <<P000361> реализация <<<P001134>,
Тип элемента <<<P000362> это
Аргумент типа <<<P000363> при <<<P001135>.

<<<P001666>{%>
Это результат сопоставления ограничений для <<<P000364>
Используя ограничение <<P000365> X<<<P000366>.
Обратите внимание, когда <<<P000367> Реализует такой класс, как <<<P001136> или <<P001137>,
не может быть подтипом <<<P001138>
(<<P000996>)%
?
<<<P001667>
Если <<P000368> реализация <<P001139>,
Выводимая пара типа ключа и значения <<P000369> является <<<P000370>,
где <<P000371> является первым и <<<P000372> Второй тип аргумента
<<<P000373> в <<<P001140>.

<<<P001668>{%>
Это результат сопоставления ограничений для <<<P000374> и <<<P000375> использовать
Ограничение <<<P000376> X<<<P000377> Y$>}$%

Обратите внимание, что этот случай и предыдущий случай
может совпадать на одном и том же элементе одновременно
когда <<P000379> Реализует оба <<P001141> и <<P001142>.
Та же ситуация возникает несколько раз ниже.
В таких случаях мы полагаемся на другие элементы, чтобы разъяснить.
?
<<<P001669>
Если <<<P000380> <<<P001670> затем
Предполагаемый тип элемента <<<P000381> является <<<P001671>,
и пара типа ключа и значения <<<P000382> это
<<P000383>.

<<<P001672>{%>
Здесь мы производим как тип элемента, так и пару типа ключа и значения.и полагаться на другие элементы для разбора.%
?
<<<P001673>
Если <<<P000384> <<<P001143> и оператор спреда <<<P001674>{...?} затем
Тип элемента <<<P000385> является <<<P001144>,
и пара типа ключа и значения <<<P000386>.
<<<P001675>
В противном случае возникает ошибка времени компиляции.
<<P000940>

<<<P001676>
% % TODO(Eernst): Уточните, почему мы не можем использовать контекстный тип
% % «Iterable» <$P_e$> Точно так же: вывод <<<P000388> в контексте «Итерируемый» <$P_e$>
%% также.
В противном случае, если <<<P000390> <<<P001145> Тогда пусть <<<P000392> быть
Предполагаемый тип <<<P000393> в контексте <<<P001146>, а затем:

<<P000911>
<<<P001677>
Если <<P000395> реализация <<P001147>,
Тип элемента <<<P000396> это
Аргумент типа <<<P000397> в <<<P001148>.
<<<P001678>{%>
Это результат сопоставления ограничений для <<<P000398> использовать
Ограничение <<<P000399> X$>}$%
?
<<<P001679>
Если <<<P000401> <<<P001680>,
Предполагаемый тип элемента <<<P000402> является <<<P001681>.
<<<P001682>
Если <<<P000403> <<<P001149> и оператор спреда <<<P001683>{...?},
Предполагаемый тип элемента <<<P000404> является <<<P001150>.
<<<P001684>
В противном случае возникает ошибка времени компиляции.
<<P000941>

<<<P001685>
В противном случае, если <<<P000405> <<<P001151> Тогда пусть <<<P000408> быть
Предполагаемый тип <<<P000409> в контексте <<<P000410>, а затем:

<<P000912>
<<<P001686>
Если <<P000411> Для реализации <<<P001152>,
Выводимая пара типа ключа и значения <<P000412> является <<<P000413>,
где <<P000414> является первым и <<<P000415> Второй тип аргументации
<<<P000416> в <<<P001153>.
<<<P001687>{%>
Это результат сопоставления ограничений для <<<P000417> и <<<P000418> использовать
Ограничение <<<P001154>.%
?
<<<P001688>
Если <<<P000422> <<<P001689>,
Выводимая пара типа ключа и значения <<<P000423> это

<<<P001690>
<<P000424>.
<<<P001691>
Если <<<P000425> <<<P001155> и оператор спреда <<<P001692>{...?},
пара типа ключа и значения <<<P000426>.
<<<P001693>
В противном случае возникает ошибка времени компиляции.
<<P000942>
<<<P001694>{-5 мм}
<<<P001695>

<<<P001339>%
<<<P001696> Если элемент
В данном случае <<<P000427> имеет форму
<<P001156> или
<<P001157>.
Условие <<P000433> Всегда выводится контекстный тип <<<P001158>.

Предположим, что "<<<P001159>" отсутствует. Тогда:

<<P000913>
<<<P001697>
Если выбранный тип элемента <<<P000435> является <<<P000436>,
Предполагаемый тип элемента <<<P000437> является <<<P000438>.
<<<P001698>
Если пара типа ключа и значения <<P000439> является <<<P000440>,
Выводимая пара типа ключа и значения <<P000441> является <<<P000442>.
<<P000943>

В противном случае Присутствует <<<P001160>.
Ошибка компиляции <<<P000444> Должен быть набор и <<P000445> Должна быть карта,
или наоборот.

<<<P001699>{%>
Это означает, что один не может распространить карту на одну ветвь и набор на другую.
С тех пор \DYNAMIC{} обеспечивает как тип элемента, так и пару типа ключа и значения,
a \DYNAMIC{} Распространение в любой ветви не приводит к возникновению ошибки.%
?

Тогда:

<<P000914>
<<<P001702>
Если выбранный тип элемента <<<P000446> является <<<P000447> и
Предполагаемый тип элемента <<<P000448> является <<<P000449>,
Тип элемента <<<P000450> это
наименьшая верхняя граница <<<P000451> и <<<P000452>.
<<<p001703>
Если выбрана пара типа ключа и значения <<<P000453> это
<<P000454>
и пара типа ключа и значения <<<P000455> это
<<P000456>,
Пара типа ключа и значения <<<P000457> это
<<P000458>,
где <<P000459> является наименьшей верхней границей <<<P000460> и <<<P000461>, и<<P000462> является наименьшей верхней границей <<<P000463> и <<<P000464>.
<<P000944>
<<<P001704>{-5 мм}
<<<P001705>

<<<P001340>>
<<<P001706> Для элемента
В данном случае <<<P000465> имеет форму
<<P001161>
где <<P000468> является производным от <<<P001707>{forLoopParts} и
<<<P001708>> Это означает, что <<<P001709> может присутствовать или отсутствовать.

Ошибки времени компиляции для <<<P000469> как
Ошибки, которые могут возникнуть при соответствующем <<<P001710> заявление
<<<P001711><<<P001712> \,\,\FOR\,\,$P$<<P002049>>>\,\{\}>,
Находится в том же объеме, что и <<P000471>.
Кроме того, ошибки и анализ типов <<<P000472> выполняется
как если бы это произошло в пределах объема тела указанного <<<P001714>>.

<<<P001715>{%>
Например, если <<P000473> имеет форму
<<P001162>
переменная <<<P001163> находится в пределах <<<P000475>.%
?

Вывод для частей
<<<<P001716>{%>
(таких, как выразительное выражение фор-ин,
или \synt{forInitializerStatement} для петли)%
?
выполнено как для соответствующего \FOR{} заявление,
включая <<<<P001719>< если и только если элемент содержит <<<P001720>.

<<P000915>
<<P001721>
Если выбранный тип элемента <<<P000476> является <<<P000477> затем
Предполагаемый тип элемента <<<P000478> является <<<P000479>.
<<<P001722>
Если пара типа ключа и значения <<P000480> является <<<P000481>,
Выводимая пара типа ключа и значения <<P000482> является <<<P000483>.
<<P000945>

<<<P001723>{%>
Другими словами, вывод течет вверх от элемента тела.
?
\vspace{3 мм)
<<<P001725>

<<<P001341>%
Наконец, мы определяем
<<<P001726>{тип вывода на наборе или карте буквально}{%
Типовой вывод! Набор или карта буквально
в целом.
Предположим, что \metavar{collection} происходит от \synt{setOrMapLiteral},
и тип контекста для \metavar{collection} является <<<P000484>.

<<P000916>
<<<P001730>
Если <<<<P001731>{коллекция} является однозначно установленным:

<<P000917>
<<P001732>
Если <<<P000485> \FreeContext затем
Статический тип <<<P001734> <<<P001164>
где <<P000487> является наименьшей верхней границей
Выводимые типы элементов.
<<P001735>
% % TODO(Eernst): Спецификация функции говорит <<<P000488>, но как мы знаем <<<P000489> Это тип?
В противном случае статический тип <<<P001736>{сборник} является <<<P000490>
где <<<P000491> определяется путем нисходящего вывода.

<<<P001737>{%>
Обратите внимание, что вывод никогда не приведет к паре типа ключа и значения.
с данным типом контекста.%
?
<<P000946>

Статический тип \metavar{сборник} является <<<P001165>.
<<P001739>
Если <<<P001740>{коллекция} является однозначной картой
где <<<P000493> <<<P001166> или <<<P000496> \FreeContext
и парами типа ключа и значения являются
<<<P001742>{K}{V}{1}{n):

Если <<<P000497> <<<P001743>{} или <<P000498> <<<P001744>,
Тип статического ключа <<<P001745> <<<P000499>
где <<P000500> является наименьшей верхней границей <<<P001746>{K}{1}{n}.
В противном случае статическим ключевым типом \metavar{collection} является $K$
где <<<P000502> определяется выводом вниз.

Если <<<P000503> <<<P001748>{} или <<P000504> <<<P001749>,
Тип статического значения \metavar{collection} $V$
где <<P000506> является наименьшей верхней границей <<<P001751>{V}{1}{n}.
В противном случае тип статического значения \metavar{collection} является $V$
где <<<P000508> определяется путем нисходящего вывода.

<<<P001753>{%>
Обратите внимание, что вывод никогда не приведет к типу элемента.
С учетом этого контекста.%
?

Статический тип <<<P001754>{сборник} является <<<P001167>.
<<P001755>В противном случае <<<P001756>{коллекция} все еще неоднозначна,
нисходящий контекст для элементов <<<P001757>
<<<P001758>>
Диагностика производится с использованием
Непосредственные элементы <<<P001759>{сборник}:

<<P000918>
<<<P001760>
Если все элементы могут быть набором,
По меньшей мере один элемент должен быть набором,
<<<<P001761>{коллекция} - набор буквальных слов с
статический тип <<<P001168>, где <<P000512> это
наименьшая верхняя граница элементов типов элементов.
<<P001762>
Если все элементы могут быть картой,
и хотя бы один элемент должен быть картой, то <<<P000513> это
Буквальная карта со статическим типом <<<P001169> <<<P000516> это
наименьшая верхняя граница ключевых типов элементов и <<P000517> это
наименьшая верхняя граница типов значений.
<<P001763>
В противном случае возникает ошибка времени компиляции.
<<<p001764>> В этом случае буквальный не может быть расшифрован. ?
<<P000947>
<<P000948>

<<<P001765>{%>
Эта последняя ошибка может произойти, если буквальный <<<P001766>{должен быть} и набором, и картой.
Вот пример: %
?

<<<P000001>

<<P001767>
<<<<P001768>{%>
Или, если ничего не указывает на то, что это <<<P001769>{или} набор или карта:%
?

<<P000002>


<\subsubsection{Сеть)
<<P001391>

<<<P001342>%
A <<<P001771>{set literal}{literal!set} обозначает заданный объект.
Правило грамматики для <<<P001772>{setOrMapLiteral}
Набор буквальных и картографических букв происходит в другом месте.
(<<P000997>).
Буквальный набор состоит из нулевых или более элементов.
(P000998).
Термин получен из <<<P001773>{setOrMapLiteral}
может быть набор буквальный или карта буквальный,
и определяется с помощью шага раздвоения
Буквально ли это набор или карта
(<<P000999>, <<<P001000>.

<<<P001343>%
При заданном наборе буквальный <<<P000518> Нет никаких аргументов,
Тип аргумента <<P000519> Выбирается так, как указано в другом месте
(<<P001001>),
и <<<P000520> отныне рассматривается как
(<<P001002>)
<<P001170>.

<<<P001774>{%>
Статический тип набора буквальных форм <<<P001171>
<<<P001172>
(<<<P001003>)%
?

<<<P001344>%
<<<P000526> быть буквальным набором формы
\code{$T$>\{\List{\ell}{1}{m}<<P002054>>>}.
Это ошибка времени компиляции, если элемент листа <<<P000528> является
<<<P001778>{mapElement}.
Это ошибка времени компиляции, если для некоторых <<P000529>
<<P000530> не имеет элемента типа,
Тип элемента <<<P000531> Не может быть назначено <<<P000532>.

<<<P001345>%
Набор может содержать ноль или более объектов.
Наборы имеют метод, который можно использовать для вставки объектов;
Это приведет к динамической ошибке, если набор не может быть изменен.
В противном случае при вставке объекта <<<P000533> в набор <<<P000534>,
Если объект <<P000535> Существует в <<<P000536> такой, что
<<P001173> Показатель <<<P001779>
не вносит никаких изменений в <<P000539>;
Если такого объекта не существует,
<<P000540> Добавлено в <<<P000541>;
В обоих случаях вставка завершается успешно.

<<<P001346>%
Набор упорядочен: итерация над элементами набора
происходит в том порядке, в котором элементы были добавлены в набор.

<<<P001780>{%>
Системные библиотеки определяют множество членов для типа <<<P001174>,
Но мы указываем только минимальный набор требований.
Используется самим языком.%

Обратите внимание, что реализация может потребовать
последовательные определения нескольких членов
класса, реализующего <<<P001175> Чтобы работать правильно.
Например, может быть геттер <<<P001176> который требуется
поведение, которое в некотором смысле согласуется с оператором; <<<P001781>{==}.
Такие ограничения документируются в системных библиотеках.
?

<<<P001347>%
Если набор буквальный <<<P000542> начинается с зарезервированного слова <<<P001782>
<<<P000543> происходит в постоянном контексте
(<<P001004>),
Это а<<<P001783>{constant set literal}{literal!set!constant}
Что является постоянным выражением
(<<P001005>)
и поэтому оценивается во время компиляции.
В противном случае это a
<<<P001784>{run-time set literal}{literal!set!run-time}
и оценивается во время выполнения.
Только буквальные значения времени выполнения могут быть мутированы после их создания.
% Эта ошибка может произойти, потому что постоянное является динамическим свойством.
Попытка мутировать постоянный набор букв приведет к динамической ошибке.

<<<P001785>{%>
% Прямо или косвенно верно следующее: Существует <<<P001786>
Модификатор % на буквальном наборе, или мы используем правило «немедленной субэкспрессии».
% — постоянные контексты.
Обратите внимание, что выражения элементов постоянного множества буквальны
происходит в постоянном контексте
(<<P001006>),
что означает <<<P001787> Модификаторы не должны быть указаны явно.%
?

<<<P001348>%
Ошибка времени компиляции, если
Литературный элемент в постоянном наборе
Это не постоянное выражение.
Ошибка времени компиляции, если
элемент в постоянном наборе буквальный
Не существует примитивного равенства.
(<<P001007>).
Ошибка времени компиляции, если два элемента постоянного набора букв равны
В соответствии с их \lit{======================================================================================================================================================================================================================================================= оператор
(<<P001008>).
Это ошибка времени компиляции, если аргумент типа постоянного множества буквален.
<<<P001789>{(независимо от того, является ли это явным или предполагаемым)}
Не является постоянным выражением типа
(<<P001009>).

<<<P001790>{%>
Связывание формального параметра типа замкнутого класса или функции
не известны во время составления,
Поэтому мы не можем использовать такие параметры типа в постоянных выражениях.
?

<<<P001349>%
Значение постоянного множества буквально
<<<P001791><<<P001792> <\,\,$T$>\{\List{\ell}{1}{m}<<< P002058>>>
является объектом <<<P000545> чей класс реализует встроенный класс
<<P001177>
где <<P000547> - фактическое значение <<P000548>
(<<<P001010>),
и чье содержание представляет собой совокупность объектов в
объектная последовательность <<<P001795>{o}{1}{n} полученный
Оценка \List{\ell}{1}{m}
(<<P001011>).
Элементы <<<P000549> происходит в том же порядке, что и
объекты в указанной последовательности объектов
<<<P001798>{(что можно наблюдать итерацией)}.

<<<P001350>%
Пусть <<<P001799>{<<<P001800>? <<<P002059><<<P002060><<<<P000550>> \{\,$\ell_{11}, \ldots, \ell_{1m_1}$\,\}
<<<P001801><<<P001802> \,\,$T_2$>\{\,$\ell_{21}, \ldots, \ell_{2m_2}$\,\}
Будьте двумя постоянными буквами.
<<<P000554> Содержание <<<P000555> Аргумент типа <<<P000556>
соответственно
<<P000557> Содержание <<<P000558> Аргумент типа <<<P000559>
быть результатом их оценки.
Далее <<<P001178> Показатель <<<P001803> если {f}
% % TODO(Eernst): См. понятие nnbd «один и тот же тип».
<<<P001179> и <<<P001180>
<<<P001804>{} для всех <<<P000566>.

<<<P001805>{%>
Другими словами, буквы с постоянным набором канонизируются, если они имеют
Один и тот же аргумент и одни и те же ценности в одном и том же порядке.
Два постоянных набора букв никогда не идентичны.
Если они имеют различное количество элементов.
Нет необходимости рассматривать канонизацию.
для других экземпляров типа <<<P001181>,
Поскольку такие случаи не
результат оценки постоянного выражения.%
?

<<<P001351>%
Набор времени выполнения \code{$T$>\{\List{\ell}{1}{n}<<P002072>>>}
Оценка проводится следующим образом:

<<P000919>
<<<P001809>
Оцениваются элементы \List{\ell}{1}{m}
(<<<P001012>),
к объектной последовательности \LiteralSequence{\List{o}{1}{n}}.
<<<P001814>
Свежий объект (<<<P001013>) <<P000568>
внедрение встроенного класса <<<P001182> Он был создан,
где <<P000570> Это действительное значение <<P000571>
(<<<P001014>).<<<P001815>
Для каждого объекта <<<P000572> в \List{o}{1}{n}, по порядку
<<P000573> Вводится в <<<P000574>.
<<<P001817>{%>
Обратите внимание, что это оставляет <<<P000575> без изменений, когда <<P000576> уже содержит
и объект <<P000577> что равно <<P000578> По словам оператора \lit{===.%
?
<<<P001819>
Результатом оценки является <<<P000579>.
<<P000949>

<<<P001352>%
Предметы, созданные установленными буквами, не переопределяются
<<<P001820>{==} оператор, унаследованный от P001183 класс.


\subsubsection{Карты}
<<P001392>

<<<P001353>%
A \IndexCustom{mapliteral}{буквально} обозначает объект карты,
Это отображение от ключей к значениям.
Правило грамматики для <<<P001823>{setOrMapLiteral}, которое охватывает оба
Буквы карты и набор букв происходит в другом месте
(<<<P001015>).
Буквальная карта состоит из нулевых или более элементов.
(<<<P001016>).
Термин получен из \synt{setOrMapLiteral}
может быть набор буквальный или карта буквальный,
и определяется с помощью шага раздвоения
Буквально ли это набор или карта
(<<<P001017>, <<<P001018>.

<<<P001354>%
При заданной карте буквальный <<<P000580> Нет никаких аргументов,
Тип аргументов <<P000581> и <<P000582> Выбирается так, как указано в другом месте
(<<<P001019>),
и <<<P000583> отныне рассматривается как
(<<<P001020>)
<<P001184>.

<<<P001825>{%>
Статический тип карты буквальный формы <<<P001185>
<<<P001186>
(\ref{setAndMapLiteralInference})%
?

<<<P001355>%
<<<P000592> Будьте картой буквальной формы
\code{$K$,\,$V$>\{\List{\ell}{1}{m}<<P002075>>>}.
Это ошибка времени компиляции, если элемент листа <<<P000595> является
<<<P001829>{выражениеЭлемент}.
Это ошибка времени компиляции, если для некоторых <<<P000596>
<<P000597> не имеет пары ключа и значения;
или парой типа ключа и значения <<P000598> является <<<P000599>,
<<<P000600> не может быть присвоено <<<P000601> или
<<P000602> Не может быть назначено <<<P000603>.

<<<P001356>%
Объект карты состоит из нуля или более записей карты.
Каждая запись имеет <<<P001040> и <<<P001041>,
Мы говорим, что карта
<<<P001830>{binds}{map!binds} или
<<<P001831>{maps}{map!maps}
Ключ к ценности.
Парой ключей и ценностей является
Добавлено на карту с помощью оператора \lit{[=],
и значение для данного ключа извлекается из карты с помощью оператора <<<P001833>{[]}.
Ключи карты обрабатываются аналогично набору
(<<<P001022>):
При связывании ключа <<<P000604> значение <<<P000605> на карте <<<P000606>
\commentary{(как в \code{<<<P000607>>>[<<<P000608>>>]\,=\,<<<P000609>>>})},
Если <<P000610> Имеется ключ <<<P000611> такой, что
<<P001188> Показатель <<<P001835>,
<<P000614> Привязывается <<<P000615> к <<<P000616>;
иначе
<<<P001836>{(когда такой ключ <<<P000617> не существует)},
Обязательный от <<<P000618> до <<<P000619> Добавлено в <<P000620>.

<<<P001357>%
Упорядочена карта: итерация по ключам, значениям или парам ключ/значение
происходит в том порядке, в котором клавиши были добавлены в набор.

<<<P001837>{%>
Системные библиотеки поддерживают множество операций на экземпляре
чей тип реализуется <<P001189>,
Но мы указываем только минимальный набор требований.
которые используются самим языком.

Обратите внимание, что реализация может потребовать
последовательные определения нескольких членов
класса, реализующего <<<P001190> Чтобы работать правильно.
Например, может быть геттер <<<P001191> который требуется
поведение, которое в некотором смысле согласуется с оператором; <<<P001838>{==}.
Такие ограничения документируются в системных библиотеках.
?

<<<P001358>%
Если карта буквальная <<P000621> начинается с зарезервированного слова <<<P001839>,
или если <<P000622> происходит в постоянном контексте
(<<<P001023>),
Это а
<<<P001840>{постоянная карта буквально}{буквально!map!constant}
Что является постоянным выражением
(<<<P001024>)и поэтому оценивается во время компиляции.
В противном случае это a
<<<<P001841>{run-time mapliteral}{literal!map!run-time}
и оценивается во время выполнения.
Только буквальные карты времени выполнения могут быть мутированы после их создания.
% Эта ошибка может произойти, потому что постоянное является динамическим свойством.
Попытка мутировать постоянную карту буквально приведет к динамической ошибке.

<<<P001842>{%>
% Прямо или косвенно верно следующее: Существует <<<P001843>
Модификатор % на буквальной карте, или мы используем правило «немедленной субэкспрессии»
% — постоянные контексты.
Обратите внимание, что ключевые и ценностные выражения постоянной карты буквальны.
происходит в постоянном контексте
(<<<P001025>),
что означает <<<P001844> Модификаторы не должны быть указаны явно.%
?

<<<P001359>%
Это ошибка времени компиляции
Литературный элемент в постоянной карте не является постоянным.
Ошибка времени компиляции, если
Ключ в постоянной карте
Не существует примитивного равенства.
(<<<P001026>).
Это ошибка времени компиляции, если два ключа постоянной карты равны.
В соответствии с их <<<P001845>{=== оператор
(<<<P001027>).
Это ошибка времени компиляции, если аргумент типа постоянной карты буквален.
<<<P001846>{(независимо от того, является ли это явным или предполагаемым)}
Не является постоянным выражением типа
(<<P001028>).

<<<P001847>{%>
Связывание формального параметра типа замкнутого класса или функции
не известны во время составления,
Поэтому мы не можем использовать такие параметры типа в постоянных выражениях.
?

<<<P001360>%
Значение постоянной карты буквально
<<<P001848><<<P001849> \,\,$T_1$\,$T_2$>\{\List{\ell}{1}{m}<<P002080>>>
является объектом <<<P000625> чей класс реализует встроенный класс
<<P001192>,
где <<<P000628> и <<<P000629> - фактическое значение <<<P000630> соответственно <<P000631>
(<<P001029>).
Пары ключей и значений <<<P000632> это
пары объектной последовательности <<<P001852>{k}{v}{1}{n} полученный
Оценка \List{<<P001854>>>}{1}{m}
(<<<P001030>),
в таком порядке.

<<<P001361>%
Пусть <<<P001855><<<P001856>? \,\,$U_1$\,$V_1$>\{\List{\ell}{1}{m_1}<<P002085>>>
<<<P001859><<<P001860> <\,<\,<$U_2$\,$V_2$>\{\List{\ell}{1}{m_2}<<<P002090>>
Будьте двумя постоянными буквами карты.
Пусть <<<P000637> с содержимым \KeyValueList{k_1}{v_1}{1}{n)
Реальные типы аргументов <<P000638>, <<P000639>
соответственно
<<P000640> с содержанием \KeyValueList{k_2}{v_2}{1}{n)
Реальный тип аргумента <<P000641>, <<P000642>
быть результатом их оценки.
Далее <<<P001193> Показатель <<<P001865> если {f}
<<<P001194>, <<<P001195>,
<<<P001196>>
<<P001197>
Для всех (P000653).

<<<P001866>{%>
Другими словами, постоянные буквальные обозначения карт канонизированы.
Нет необходимости рассматривать канонизацию.
Для других типов <<P001198>,
Поскольку такие случаи не могут быть
результат оценки постоянного выражения.%
?

<<<P001362>%
Карта времени выполнения буквальная
<\code{$T_1, T_2$>\{\List{\ell}{1}{m}<<P002092>>>}}
Оценка проводится следующим образом:

<<P000920>
<<<P001870>
Оцениваются элементы \List{\ell}{1}{m}
(<<<P001031>),
в объектную последовательность \LiteralSequence{\KeyValueList{k}{v}{1}{n}}.
<<<P001875>
Свежий пример (<<<P001032>) <<P000655>
чей класс реализует встроенный класс <<<P001199>
выделяется,
где <<<P000658> и <<<P000659> являются действительными значениями <<<P000660> соответственно <<P000661>
(<<P001033>).
<<<P001876>
Оператор \lit{[]=} вызывается на $o$
с первым аргументом <<P000663> и второй аргумент <<<P000664>,
Для каждого <<<P000665>, в таком порядке.
<<<P001878>Результатом оценки является <<P000666>.
<<P000950>

<<<P001363>%
Объекты, созданные буквальными картами, не переопределяются
тот <<<P001879>{=== оператор, унаследованный от <<<P001200> класс.


<<<P001880> Бросок.
<<P001393>

<<<P001364>%
<<<P001042> Используется для исключения.

<<P000921>
<rowExpression> ::= <<<<P001881>< <выражение>

<rowExpressionWithoutCascade> ::= \THROW{} <expressionWithoutCascade>
<<P000951>

<<<P001365>%
Оценка броскового выражения формы
\code{<<P001884>>>{} <<P000667>;
идет следующим образом:

<<<P001366>%
Выражение <<P000668> оценивается на предмет <<P000669>
(<<<P001034>).

<<<P001885>{%>
Нет требования, чтобы выражение <<P000670> должны оценивать
Любые специальные объекты.%
?

<<<P001367>%
Если <<P000671> является нулевым объектом (<<<P001035>), затем бросается <<<P001201>.
В противном случае пусть <<P000672> быть следом стека, соответствующим текущему состоянию исполнения,
и <<<P001886> Заявление с броском <<<P000673> как объект исключения
<<P000674> След стека (<<<P001036>).

<<<P001368>%
Если <<P000675> является примером класса <<<P001202> или его подкласса,
Это первый случай, когда <<<P001203> объект выбрасывается,
След стека <<<P000676> Сохраняется на <<<P000677> Чтобы он был возвращен
<<<P001204> Геттер унаследован от <<<P001205>.

<<<P001887>{%>
Если то же самое <<<P001206> объект выбрасывается более одного раза,
<<<P001207> Геттер вернет стековый след
<<<<P001888 >> {первое время было выброшено.%
?

<<<P001369>%
Статический тип выражения броска <<<P000678>.


<<<<P001889>{Функциональные выражения}
<<P001394>

<<<P001370>%
A <<<P001890>{функция дословная}{дословная!функция}
является анонимной декларацией и выражением
который инкапсулирует исполняемый блок кода.

% % TODO(Eernst): Это очень неоднозначно, потому что <первичный> получает
% <выражение функции>. Dart.g извлекает <functionExpression> из
% % <expression>, но допускает только блок в <functionPrimary>, который
% является производным от <primary>. Функция «ExpressionWithoutCascade»
%%, что позволяет выполнять функцию «=» буквально в каскадном задании.
% % Тем не менее, мы должны очень тщательно оценить поломку, прежде чем принимать
% от этого подхода, поскольку он предотвращает разбор буквальных функций
%% как условное выражение, если нулевое выражение, ... постфиксное выражение.
<<P000922>
<functionExpression> ::= <formalParameterPart> <functionExpressionBody>

<functionExpressionBody> ::= <<<P001891>? <=> <выражение>
<<<P001892> (\ASYNC{} '*'? | \SYNC{} '*')) <блокировать>
<<P000952>

<<<P001371>%
Грамматика не позволяет функции буквально объявить тип возврата,
Но возможно, что функция, дословно имеющая
<<<<P001895>{объявлено о возвращении Тип {буквальный!функциональный!декларированный тип возврата},
потому что она может быть получена посредством вывода типа.
Такой тип возврата включает
Когда мы говорим об объявленном типе функции возврата.

<<<P001896>{%>
Тип вывода будет указан в будущей версии этого документа.
В настоящее время мы рассматриваем вывод типа как завершенный этап,
Этот документ определяет значение программ Дарта.
где уже добавлены предполагаемые типы.%
?

<<<P001372>%
Мы говорим, что тип <<<P000679>
<<<P001897>{derives a future type}{type!derives a future type}
<<P000680> в следующих случаях, используя первый применимый случай:

% % TODO(Eernst): Обратите внимание, что «X расширяет X?» может создать бесконечный цикл.
% % Возможно, мы сделаем такую ошибку.
<<P000923>
<<<P001898>
% % TODO(Eernst): Приходите классы микширования и типы расширений: добавьте их.
Если <<P000681> Это тип, который вводится
класс, микширование или декларирование enum,
Если <<P000682> Прямой или косвенный суперинтерфейс
(<<<P001037>)
<<<P000683> <<<P001208> для некоторых <<<P000685>,<<<P000686> Получает будущий тип <<<P001209>.
<<<P001899>
Если <<P000688> Тип <<<P001210> для некоторых <<<P000690>,
<<<P000691> Получает будущий тип <<<P001211>.
<<<p001900>
Если <<<P000693> <<<P001212> Для некоторых <<<P000695>
<<P000696> Получает будущий тип <<P000697>,
<<<P000698> Получает будущий тип <<<P001213>.
<<<p001901>
Если <<P000700> является переменной типа со связанным <<<P000701>, и
<<P000702> Получает будущий тип <<<P000703>,
<<<P000704> Получает будущий тип <<P000705>.
<<<p001902>
<<<P001903>{%>
Не существует правила для случая, когда <<P000706> имеет форму <<<P001214>
Потому что этого никогда не произойдет
(это понятие используется только в <<<P001904>, который определен ниже.)
?
<<P000953>

<<<P001373>%
Если ни один из этих случаев не применим,
Мы говорим, что <<P000709> Он не имеет будущего типа.

<<<P001905>{%>
Обратите внимание, что если <<<P000710> Получает будущий тип <<<P000711> <<<<P001906>> T}{F},
<<P000712> всегда имеет форму <<<P001215> или <<<P001216>,
где <<<P000715> <<<P001217> или \code{FutureOr}. Доказательство является
по индукции на структуру <<<P000716>:

<<P000924>
<<<p001907>
% % TODO(Eernst): Приходите классы микширования и типы расширений: добавьте их.
Если <<P000717> Это тип, который вводится
класс, микширование или декларирование enum,
Если <<P000718> Прямой или косвенный суперинтерфейс
(<<<P001038>)
<<<P000719> <<<P001219> Для некоторых <<<P000721>, тогда
<<<P001220> и <<P001221>, <<P000725>.
<<<p001908>
Если <<P000726> Тип <<<P001222> Для некоторых <<<P000728>, то по рефлексивности,
<<<P001223>. <<<P001224> и
<<<P001225>, следует, что <<P000734>.
<<<p001909>
Если <<<P000735> <<<P001226> Для некоторых <<<P000737>
<<P000738> Получает будущий тип <<<P000739>,
С помощью индукционной гипотезы, <<<P001227>,
где <<P000741> имеет форму
<<<P001228> или <<<P001229> и <<P000746> <<<P001230> или
<<<P001231>. Следовательно, <<<P001232>, и путем замены,
<<<P001233>. С тех пор <<<P001234> для всех <<<P000752>
<<<P001235>. Итак, позволив <<<P000755> и <<<P001236>,
Отсюда следует, что <<P000758>.
<<<P001910>
Если <<P000759> является переменной типа со связанным <<<P000760>, и
<<P000761> Получает будущий тип <<P000762>,
С помощью индукционной гипотезы, <<<P001237>,
где <<P000764> имеет форму
<<<P001238> или <<P001239> и <<P000769> <<<P001240> или
<<P001241>. Кроме того, с тех пор <<<P000770> является пределом <<<P000771>, <<<P000772>>
транзитивность, <<P000773>. Таким образом, позволяя <<<P000774> и <<<P000775>, это
Далее следует <<P000776>.
<<P000954>

Также обратите внимание, что «производные» в этом контексте относятся к вычислениям.
где тип <<<P000777> Приведены супертипы <<<P000778> обыскиваются,
и тип <<<P000779> Выбирается одна из этих форм.
Нет никакой связи с понятием «производный класс», означающим «подкласс». "
Некоторые языковые сообщества используют.%
?

<<<P001374>%
Определяем вспомогательную функцию
\IndexCustom{<<P001912>>>{T}}{flatten(t)@<<P001913>>>{flatten}<<P000780>>>}
следующим образом, используя первый применимый случай:

<<P000925>
<<<P001914> Если <<<P000781> <<<P001242>
переменная типа <<<P000784> Тип <<P000785> затем

<<P000926>
<<<P001915>, если <<P000786> Получает будущий тип <<P000787>
затем <\DefEquals{\flatten{T}}{\code{\flatten !
<<<P001920> В противном случае
\DefEquals{<<P001922>>>{T}}{<<P001923>>>{<<P001924> Икс.
<<P000955>

<<<P001925> Если <<P000788> Получает будущий тип <<P001243>
<<<P001244><<<P001926><<<P001927> T}}{S}

<<<P001928> Если <<P000791> Получает будущий тип <<<P001245><<<P002093>
<<<P001246>\  \DefEquals{\flatten{T}}{\code{<<<P000794>>>?}}.

<<<p001931> В противном случае \DefEquals{\flatten{ T}}{T}.
<<P000956>

<<<P001934>{%>
Это определение гарантирует, что для любого типа <<<P000795>
<<P001248>. Доказательством является индукция на
Структура <<<P000798>:

<<P000927>
<<<p001935> Если <<<P000799> <<<P001249> затем

<<P000928>
<<<P001936>, если <<P000802> Получает будущий тип <<<P000803>,
Затем <-<<P001250> и <<<P001251>, так <<<P001252>.
По индукционной гипотезе <<<P001253>.
С тех пор <<<P001254> В данном случае следует, что
<<<<P001255>>
<<<P001256>.
<<<p001937> В противном случае <<P001257>.
По индукционной гипотезе, <<<P001258>.
С тех пор <<P001259> В данном случае следует, что
<<<<P001260>>
<<P001261>.
<<P000957>

<<<P001938> Если <<P000822> Получает будущий тип <<P001262>
или <<<P001263>, то, поскольку <<<P001264>,
Отсюда следует, что <<<P001265>. С тех пор <<P001266>
Из этого следует, что <<<P001267>.

<<<P001939> Если <<P000833> Получает будущий тип <<P001268> или
<<<P001269>, с тех пор <<<P001270>,
Отсюда следует, что <<<P001271>.
<<P001272> для любого типа <<P000844>
(это может быть показано с использованием правил подтипа союзного типа и
<<<P001273> по ковариантности, т. е. по транзитивности,
<<P001274>. С тех пор <<P001275> в данном случае,
Отсюда следует, что <<<P001276>.

<<<P001940> В противном случае, <<<P001277>
<<<P001278>. С тех пор
<<<<P001279>> Из этого следует, что
<<P001280>.
<<P000958>
?

<<<P001375>%
<<<P001941> Позиция, стрелка.
Статический тип функции буквальный формы

<<<P001942>
<<P001281>

<<<P001943>
<<P001282>

<<P001944>
это
<<<P001945>{T_0}

<<<p001946>
% % [вывод]: Статический тип функции буквально может быть выведен.
где <<P000864> Статический тип <<<P000865>.
<<P001947>

<<<P001376>%
<<<P001948> Позиция, стрелка, будущее.
Статический тип функции буквальный формы

<<<P001949>
<<<P001283>

<<<P001950>
<<<P001951> (<<P000866>) [<<P000867>] \ASYNC{} => $e$>

<<<P001953>
это
<<<P001954><<<P001955> Будущее<\flatten{T_0}>

<<<P001957>
где <<P000869> Статический тип <<<P000870>.
<<<P001958>

<<<P001377>%
<<<P001959> Названа, стрелка.
Статический тип функции буквальный формы

<<<P001960>
<<P001284>

<<<P001961>
<<<P001962><<<P000871> \{$T_{n+1}\ x_{n+1} = d_1, \ldots,\ T_{n+k}\ x_{n+k} = d_k$\}> <<P000873>

<<<p001963>
это
<<<P001964>{T_0}

<<<P001965>
где <<P000874> Статический тип <<<P000875>.
<<<p001966>

<<<P001378>%
<<<P001967>{Название, стрелка, будущее}
Статический тип функции буквальный формы

<<<P001968>
<<P001285>

<<<P001969>
<<<P001970> (<<P000876>) \{$T_{n+1}\ x_{n+1} = d_1, \ldots,\ T_{n+k}\ x_{n+k} = d_k$\} \ASYNC{} => $e$

<<<P001972>
это
<<<P001973><<<P001974>> Будущее<\flatten{T_0}>

<<<P001976>
где <<P000879> Статический тип <<<P000880>.
<<P001977>

<<<P001379>%
<<<P001978> Позиционно, блок.
Статический тип функции буквальный формы

<<<p001979>
\code{\TypeParameters{X}{B}{S}>

<<<P001982>
<<<P001983> (<<P000881>) [<<P000882>] <<<P002099> <<P000883> <<<P002100>>

<<<P001984>
это
<<<P001985><<<P001986>>
<<<p001987>

<<<P001380>%
<<<P001988> Позиция, блок, будущее.
Статический тип функции буквальный формы

<<<P001989>
<<P001286>

<<<P001990>
<<<P001991> (<<P000884>) [<<P000885>] <<<P001992>< <<<p002101> <<P000886> <<<P002102>>

<<<p001993>
это
% % TODO(Eernst): Настройка для учета выводов типа.<<<P001994><<<P001287>>.
<<<P001995>

<<<P001381>%
<<<p001996> Позиция, блок, поток.
Статический тип функции буквальный формы

<<<P001997>
<<P001288>

<<<P001998>
<<<P001999> (<<P000887>) [<<P000888>] <<<P002000>< <<<P002103> <<P000889> <<<P002104>>

<<<P002001>
это
% % TODO(Eernst): Настройка для учета выводов типа.
\FunctionTypePositionalStdCr{\code{Stream}}.
<<P002003>

<<<P001382>%
<<<P002004>{Позиционное, блоковое, итерабельное}
Статический тип функции буквальный формы

<<P002005>
<<P001290>

<<<P002006>
<<<P002007> (<<P000890>) [<<P000891>] <<<P002008>< <<<P002105> <<P000892> <<<P002106>>

<<<P002009>
это
% % TODO(Eernst): Настройка для учета выводов типа.
<<<P002010><<<P001291>>.
<<<P002011>

<<<P001383>%
<<<<P002012>{Название, блок)
Статический тип функции буквальный формы

<<<p002013>>
<<P001292>

<<<P002014>
<<<P002015> (<<P000893>) [<<P000894>] <<<p002107> <<P000895> <<<P002108>>

<<<P002016>
это
% % TODO(Eernst): Настройка для учета выводов типа.
<<<P002017><<<P002018>>.
<<<P002019>

<<<P001384>%
<<<P002020>{Название, блок, будущее}
Статический тип функции буквальный формы

<<<P002021>
<<P001293>

<<<P002022>
<<<P002023> (<<P000896>) \{$T_{n+1}\ x_{n+1} = d_1, \ldots,\ T_{n+k}\ x_{n+k} = d_k$\} <<<<P002024 >> <<P002111> $s$\}

